% GENERATED FILE. Edit canonical lessons, then run npm run generate:books.
% canonical-source-hash: fnv1a64:e28e72565e37116a
% canonical-lessons: GU-C269-sambhalvu, GU-C269-uchervu, GU-C269-vakhanvu, GU-C269-yad-avvi, GU-C269-adar

\chapter{Caring for each other}
\label{ch:gu-caring-for-each-other}
\hlchaptermodality{269}

\begin{chapteropening}
\textbf{By the end of this chapter:} \emph{I can ask who looks after whom and say whom I miss.}

Complete the last lesson of chapter 269: I can ask who looks after whom and say whom I miss.
\end{chapteropening}

\section[\texorpdfstring{\gu{સંભાળવું} (sambhāḷvũ)}{સંભાળવું (sambhāḷvũ)}]{\gu{સંભાળવું} (sambhāḷvũ) — to look after, to take care of}

\label{lesson:GU-C269-sambhalvu}



\begin{quote}
Before the new one: say the Gujarati for to persuade, then the Gujarati for advice.
\end{quote}

\subsection*{You'll want to know: \gu{સંભાળવું}}
\textbf{\gu{સંભાળવું}} — \emph{sambhāḷvũ} — \textquotedblleft{}to look after, to take care of\textquotedblright{}.

\begin{quote}
\textbf{\gu{શું} \gu{તમે} \gu{તમારી} \gu{દાદીને} \gu{સંભાળો} \gu{છો}?} — \emph{śũ tame tamārī dādīne sambhāḷo cho?} — \textquotedblleft{}Do you look after your grandmother?\textquotedblright{}
\end{quote}

Not the taught \textbf{\gu{સાંભળવું}}, \emph{sā̃bhaḷvũ}, \textquotedblleft{}to hear\textquotedblright{}.

\begin{cousinweb}
\textbf{\gu{સંભાળવું}} (\emph{sambhāḷvũ}) is a tadbhava. CDIAL lists it, glossed \textquotedblleft{}to guard, take care of\textquotedblright{}, under Sanskrit \emph{sáṁbharati}, \textquotedblleft{}brings together\textquotedblright{}, and its causative \emph{saṁbhārayati}, \textquotedblleft{}maintains\textquotedblright{} (CDIAL 12961).
\end{cousinweb}

\begin{grammarlens}[title={What you've built}]
One more word for everyday life at A2.
\end{grammarlens}

\subsection*{Your turn}
\begin{itemize}
\raggedright
  \item \emph{Say it:} \emph{sambhāḷvũ}
  \item \emph{Say it:} \emph{sambhāḷvũ}, once more
  \item \emph{Say it:} \emph{śũ tame tamārī dādīne sambhāḷo cho?} — Do you look after your grandmother?
  \item \emph{Recall:} say the Gujarati for to persuade, then the Gujarati for advice, then say \emph{sambhāḷvũ} again
\end{itemize}

\begin{tcolorbox}[breakable,colback=teal!4,colframe=teal!35!black,title={Before you move on}]
What does \gu{સંભાળવું} mean? (\textquotedblleft{}To look after, to take care of\textquotedblright{}.) Say it once more.
\end{tcolorbox}
\section[\texorpdfstring{\gu{ઉછેરવું} (uchervũ)}{ઉછેરવું (uchervũ)}]{\gu{ઉછેરવું} (uchervũ) — to bring up, to raise}

\label{lesson:GU-C269-uchervu}



\begin{quote}
Before the new one: say the Gujarati for advice, then the Gujarati for to look after.
\end{quote}

\subsection*{You'll want to know: \gu{ઉછેરવું}}
\textbf{\gu{ઉછેરવું}} — \emph{uchervũ} — \textquotedblleft{}to bring up, to raise\textquotedblright{}.

\begin{quote}
\textbf{\gu{શું} \gu{તમારી} \gu{દાદીએ} \gu{તમને} \gu{ઉછેર્યા}?} — \emph{śũ tamārī dādīe tamne ucheryā?} — \textquotedblleft{}Did your grandmother bring you up?\textquotedblright{}
\end{quote}

\begin{cousinweb}
\textbf{\gu{ઉછેરવું}} (\emph{uchervũ}) is inherited, the transitive partner of \emph{ucharvũ}, \textquotedblleft{}to be brought up\textquotedblright{}. CDIAL derives \emph{ucharvũ} from \emph{*ut-śrīyatē}, Sanskrit \emph{úcchrīyatē}, \textquotedblleft{}is raised\textquotedblright{} (CDIAL 1864).
\end{cousinweb}

\begin{grammarlens}[title={What you've built}]
One more word for everyday life at A2.
\end{grammarlens}

\subsection*{Your turn}
\begin{itemize}
\raggedright
  \item \emph{Say it:} \emph{uchervũ}
  \item \emph{Say it:} \emph{uchervũ}, once more
  \item \emph{Say it:} \emph{śũ tamārī dādīe tamne ucheryā?} — Did your grandmother bring you up?
  \item \emph{Recall:} say the Gujarati for advice, then the Gujarati for to look after, then say \emph{uchervũ} again
\end{itemize}

\begin{tcolorbox}[breakable,colback=teal!4,colframe=teal!35!black,title={Before you move on}]
What does \gu{ઉછેરવું} mean? (\textquotedblleft{}To bring up, to raise\textquotedblright{}.) Say it once more.
\end{tcolorbox}
\section[\texorpdfstring{\gu{વખાણવું} (vakhāṇvũ)}{વખાણવું (vakhāṇvũ)}]{\gu{વખાણવું} (vakhāṇvũ) — to praise}

\label{lesson:GU-C269-vakhanvu}



\begin{quote}
Before the new one: say the Gujarati for to look after, then the Gujarati for to bring up.
\end{quote}

\subsection*{You'll want to know: \gu{વખાણવું}}
\textbf{\gu{વખાણવું}} — \emph{vakhāṇvũ} — \textquotedblleft{}to praise\textquotedblright{}.

\begin{quote}
\textbf{\gu{સાહેબ} \gu{ક્યારેય} \gu{કોઈને} \gu{નથી} \gu{વખાણતા}.} — \emph{sāheb kyārey koīne nathī vakhāṇtā} — \textquotedblleft{}The boss never praises anyone.\textquotedblright{}
\end{quote}

\begin{cousinweb}
\textbf{\gu{વખાણવું}} (\emph{vakhāṇvũ}) is a tadbhava. CDIAL derives it from Sanskrit \emph{vyākhy\'{\={a}}na}, \textquotedblleft{}comment, narration\textquotedblright{}, by way of the noun \emph{vakhāṇ}, \textquotedblleft{}praise\textquotedblright{} (CDIAL 12188).
\end{cousinweb}

\begin{grammarlens}[title={What you've built}]
One more word for everyday life at A2.
\end{grammarlens}

\subsection*{Your turn}
\begin{itemize}
\raggedright
  \item \emph{Say it:} \emph{vakhāṇvũ}
  \item \emph{Say it:} \emph{vakhāṇvũ}, once more
  \item \emph{Say it:} \emph{sāheb kyārey koīne nathī vakhāṇtā} — The boss never praises anyone.
  \item \emph{Recall:} say the Gujarati for to look after, then the Gujarati for to bring up, then say \emph{vakhāṇvũ} again
\end{itemize}

\begin{tcolorbox}[breakable,colback=teal!4,colframe=teal!35!black,title={Before you move on}]
What does \gu{વખાણવું} mean? (\textquotedblleft{}To praise\textquotedblright{}.) Say it once more.
\end{tcolorbox}
\section[\texorpdfstring{\gu{યાદ} \gu{આવવી} (yād āvvī)}{યાદ આવવી (yād āvvī)}]{\gu{યાદ} \gu{આવવી} (yād āvvī) — to remember fondly, to miss}

\label{lesson:GU-C269-yad-avvi}



\begin{quote}
Before the new one: say the Gujarati for to bring up, then the Gujarati for to praise.
\end{quote}

\subsection*{You'll want to know: \gu{યાદ} \gu{આવવી}}
\textbf{\gu{યાદ} \gu{આવવી}} — \emph{yād āvvī} — \textquotedblleft{}to remember fondly, to miss\textquotedblright{}.

\begin{quote}
\textbf{\gu{શું} \gu{તમને} \gu{ઘરની} \gu{યાદ} \gu{આવે} \gu{છે}?} — \emph{śũ tamne gharnī yād āve che?} — \textquotedblleft{}Do you miss home?\textquotedblright{}
\end{quote}

The one who misses takes \textbf{-\gu{ને}}: \textbf{\gu{મને} \gu{તારી} \gu{યાદ} \gu{આવે} \gu{છે}}, \emph{mane tārī yād āve che}, \textquotedblleft{}I miss you\textquotedblright{}.

\begin{cousinweb}
\textbf{\gu{યાદ} \gu{આવવી}} (\emph{yād āvvī}), for a memory to come, is built from \emph{yād}, \textquotedblleft{}memory\textquotedblright{}, a loan usually traced to Persian \emph{yād}, and the verb \emph{āvvī}, \textquotedblleft{}to come\textquotedblright{}.
\end{cousinweb}

\begin{grammarlens}[title={What you've built}]
One more word for everyday life at A2.
\end{grammarlens}

\subsection*{Your turn}
\begin{itemize}
\raggedright
  \item \emph{Say it:} \emph{yād āvvī}
  \item \emph{Say it:} \emph{yād āvvī}, once more
  \item \emph{Say it:} \emph{śũ tamne gharnī yād āve che?} — Do you miss home?
  \item \emph{Recall:} say the Gujarati for to bring up, then the Gujarati for to praise, then say \emph{yād āvvī} again
\end{itemize}

\begin{tcolorbox}[breakable,colback=teal!4,colframe=teal!35!black,title={Before you move on}]
What does \gu{યાદ} \gu{આવવી} mean? (\textquotedblleft{}To remember fondly, to miss\textquotedblright{}.) Say it once more.
\end{tcolorbox}
\section[\texorpdfstring{\gu{આદર} (ādar)}{આદર (ādar)}]{\gu{આદર} (ādar) — respect}

\label{lesson:GU-C269-adar}



\begin{quote}
Before the new one: say the Gujarati for to praise, then the Gujarati for to remember fondly.
\end{quote}

\subsection*{You'll want to know: \gu{આદર}}
\textbf{\gu{આદર}} — \emph{ādar} — \textquotedblleft{}respect\textquotedblright{}.

\begin{quote}
\textbf{\gu{શું} \gu{તમે} \gu{તમારા} \gu{ગુરુનો} \gu{આદર} \gu{કરો} \gu{છો}?} — \emph{śũ tame tamārā guruno ādar karo cho?} — \textquotedblleft{}Do you respect your teacher?\textquotedblright{}
\end{quote}

\begin{cousinweb}
\textbf{\gu{આદર}} (\emph{ādar}) is Sanskrit \emph{ādara}, \textquotedblleft{}respect\textquotedblright{}, taken over whole as a tatsama (CDIAL 1142); EtymDB also lists it as borrowed from Sanskrit.
\end{cousinweb}

\begin{grammarlens}[title={What you've built}]
One more word for everyday life at A2.
\end{grammarlens}

\subsection*{Your turn}
\begin{itemize}
\raggedright
  \item \emph{Say it:} \emph{ādar}
  \item \emph{Say it:} \emph{ādar}, once more
  \item \emph{Say it:} \emph{śũ tame tamārā guruno ādar karo cho?} — Do you respect your teacher?
  \item \emph{Recall:} say the Gujarati for to praise, then the Gujarati for to remember fondly, then say \emph{ādar} again
\end{itemize}

\begin{tcolorbox}[breakable,colback=teal!4,colframe=teal!35!black,title={Before you move on}]
What does \gu{આદર} mean? (\textquotedblleft{}Respect\textquotedblright{}.) Say it once more.
\end{tcolorbox}
